\notechapter{N-20220803}{Arithmetic and Geometric Sequence Exercises}

\section{A trigonometric opening example}

The first exercise concerns a triangle \(ABC\) such that
\[
  \lg(\sin A),\quad \lg(\sin B),\quad \lg(\sin C)
\]
form an arithmetic progression, and \(A,B,C\) themselves also form an arithmetic progression.  Since
\[
  A+B+C=\pi,
\]
if \(A,B,C\) are in arithmetic progression, then
\[
  B=\frac\pi3.
\]
The logarithmic condition means
\[
  2\lg(\sin B)=\lg(\sin A)+\lg(\sin C),
\]
so
\[
  \sin^2B=\sin A\sin C.
\]
Because \(\sin B=\sqrt3/2\), we get
\[
  \sin A\sin C=\frac34.
\]
Using \(A+C=2\pi/3\), one can analyze
\[
  \cos(A-C)=\cos A\cos C+\sin A\sin C.
\]
The computation in the page leads to the conclusion that the triangle is equilateral.  The underlying idea is that two independent arithmetic-progressions conditions force symmetry.

\section{Logarithmic sequences}

If
\[
  a_n=10+\lg 2^n,
\]
then
\[
  a_n=10+n\lg2.
\]
Hence
\[
  a_{n+1}-a_n=\lg2,
\]
so \(\{a_n\}\) is an arithmetic progression.

This is a typical use of logarithms: they turn powers into multiples.  A geometric pattern in \(2^n\) becomes an arithmetic pattern after applying \(\lg\).

\section{Recovering terms from partial sums}

If \(S_n\) is the sum of the first \(n\) terms of a sequence, then
\[
  a_n=S_n-S_{n-1}.
\]
This applies to problems involving two arithmetic progressions \(\{a_n\}\), \(\{b_n\}\) whose sums \(S_n,T_n\) satisfy
\[
  \frac{S_n}{T_n}=\frac{2n}{3n+1}.
\]
For arithmetic progressions,
\[
  S_n=\frac n2(a_1+a_n),
\]
but the cleaner trick is to use
\[
  a_n=S_n-S_{n-1},\qquad b_n=T_n-T_{n-1}.
\]
The page obtains
\[
  \frac{a_n}{b_n}=\frac{4n-2}{6n-2}=\frac{2n-1}{3n-1}.
\]

\section{Sums and recurrence relations}

A common type of problem gives a relation between \(S_n\) and \(a_n\).  For example,
\[
  S_n=2a_n-2.
\]
Then
\[
  a_{n+1}=S_{n+1}-S_n=2a_{n+1}-2-(2a_n-2),
\]
so
\[
  a_{n+1}=2a_n.
\]
Thus the sequence is geometric.  This is a useful principle: whenever a problem relates a term to a partial sum, subtract consecutive equations.

\section{Geometric progression products}

If \(\{a_n\}\) is a geometric progression, then
\[
  a_m^2=a_{m-r}a_{m+r}.
\]
This appears in several exercises.  It follows from
\[
  a_n=a_1q^{n-1}.
\]
For example,
\[
  a_6^2=a_5a_7.
\]
Such identities compute missing terms and ratios from partial information.

\section{A harmonic-looking sequence}

One page studies
\[
  a_n=\frac1{1+2}+\frac1{2+3}+\cdots+\frac1{n+n+1}
\]
or similar telescoping sums.  The key transformation is
\[
  \frac1{k(k+1)}=\frac1k-\frac1{k+1}.
\]
Telescoping is the additive analogue of cancellation in products.  When the denominator factors into consecutive terms, always look for a difference of reciprocals.

\section{A final synthesis}

The underlying methods are stable across the exercises: logarithms convert multiplication into addition, partial sums are inverted by differences, arithmetic progressions have linear terms, geometric progressions have exponential terms, and telescoping makes hidden cancellation explicit.  These are the basic patterns of finite-difference calculus.

\section{Sequences as discrete dynamical systems}

An arithmetic progression and a geometric progression are both generated by repeatedly applying a simple transformation.  For an arithmetic progression,
\[
  a_{n+1}=a_n+d,
\]
so the update is translation.  For a geometric progression,
\[
  a_{n+1}=qa_n,
\]
so the update is scaling.  This explains why logarithms turn geometric progressions into arithmetic ones:
\[
  \log a_{n+1}=\log q+\log a_n.
\]
Multiplicative iteration becomes additive iteration after applying \(\log\).

This is the same structural idea behind linear recurrences.  A sequence is not only a list of numbers; it is often an orbit under an update rule.

\section{Partial sums and finite differences}

If \(S_n=a_1+\cdots+a_n\), then
\[
  a_n=S_n-S_{n-1}=\Delta S_{n-1}.
\]
So recovering terms from partial sums is finite differentiation.  Conversely,
\[
  S_n=S_0+\sum_{k=1}^{n}a_k
\]
is finite integration.

For an arithmetic progression \(a_n=a_1+(n-1)d\), the first difference is constant:
\[
  \Delta a_n=d.
\]
For a quadratic sequence, the second difference is constant.  This is why finite-difference tables detect polynomial sequences, just as derivatives detect polynomial degree in calculus.

\section{Generating functions for sequence exercises}

A sequence \((a_n)\) can be packaged into a generating function
\[
  A(x)=\sum_{n\ge0}a_nx^n.
\]
Arithmetic and geometric progressions have especially simple forms.  For example,
\[
  \sum_{n\ge0}q^nx^n=\frac{1}{1-qx}.
\]
A linear sequence produces a rational generating function with denominator \((1-x)^2\).  Recurrences translate into algebraic equations for generating functions.

This turns many sequence problems into algebra: instead of manipulating terms one by one, encode the whole sequence and operate on the generating function.

\section{Telescoping as a coboundary}

A telescoping sum has the form
\[
  a_k=F(k+1)-F(k).
\]
Then
\[
  \sum_{k=m}^{n}a_k=F(n+1)-F(m).
\]
Only the boundary terms survive.  In this language, a telescoping term is a discrete derivative or coboundary.

The identity
\[
  \frac{1}{k(k+1)}=\frac1k-\frac{1}{k+1}
\]
is exactly this pattern with \(F(k)=1/k\).  The cancellation is not accidental; it is the discrete fundamental theorem of calculus.

\section{Second-order recurrences and characteristic roots}

Some sequence problems go beyond arithmetic and geometric progressions.  A linear recurrence such as
\[
  a_{n+2}=pa_{n+1}+qa_n
\]
is solved by trying \(a_n=\lambda^n\).  This gives the characteristic equation
\[
  \lambda^2=p\lambda+q.
\]
When the roots are distinct, the solution has the form
\[
  a_n=A\lambda_1^n+B\lambda_2^n.
\]
Thus geometric progressions are the eigenvectors of the shift recurrence.  This is the bridge from elementary sequences to linear algebra and dynamical systems.

\section{Discrete calculus of elementary sequences}

The exercises are unified by discrete calculus.  Logarithms linearize multiplication, finite differences invert partial sums, telescoping records boundary cancellation, generating functions encode whole sequences, and linear recurrences are controlled by characteristic roots.  The elementary tricks are the visible form of a general principle: choose the representation in which the sequence's update rule becomes simple.
